\begin{theorem}[Nash--Williams]
Let $F$ be a front.  For every subset $S$ of $F$ there is a sub-front
$F'\subseteq F$ such that either $F'\subseteq S$ or
$F'\cap S=\emptyset$.
\end{theorem}
\begin{proof}
We prove the assertion by transfinite induction on the ordinal rank of the
prefix tree.  This induction is legitimate because the relation in
\noderef{FrontTreeWellFounded} is well-founded, and its orientation is
important: a relation step from a node $v$ to a node $u$ means that $u$ is a
proper extension of $v$, since it is
\(\mathsf{ProperInitialSegment}(v,u)\) in the notation of
\noderef{ProperInitialSegment}.  Thus the induction hypothesis is available
for every strict successor subtree.

Here is the precise rank comparison used below.  If $F$ is nontrivial on
$X$ and $n\in X$, the map
\[
 \phi_n(s)=\{n\}\cup s
\]
identifies the prefix tree of the ray $F_n$ with
\[
 \{u\in\mathsf{PrefixTree}(F)\mid \{n\}\segm u\},
\]
the successor subtree rooted at the strict successor $\{n\}$ of the empty
root.  Indeed, if
$s\in\mathsf{PrefixTree}(F_n)$, then by the definitions of
\noderef{PrefixTree} and \noderef{Ray}, $s$ is an initial segment of a
member of $F_n$, and inserting the common least element $n$ makes
$\phi_n(s)$ an initial segment of the corresponding member of $F$.
Conversely, if $u$ is in the displayed successor subtree, then
\noderef{PrefixTree} gives $t\in F$ with $u\segm t$; the initial-segment
description in \noderef{InitialSegment} shows that both $u$ and $t$ begin
with $n$, and deleting that common first element gives a prefix of the ray
member $t\setminus\{n\}$ in \noderef{Ray}.  This proves the claimed
identification and preserves proper extension.  Consequently the rank of
$F_n$ is strictly smaller than the rank of $F$; explicitly, the paper's
ray/rank formula is
\[
 \rk F=\sup\{\rk(F_n)+1\mid n\in X\}.
\]
Thus the induction hypothesis applies to every $F_n$.

The trivial case is immediate.  If $F=\{\emptyset\}$, then
$S\subseteq F$ is either empty or all of $F$.  Take $F'=F$ and $Y=X$;
the front assertion is the hypothesis from \noderef{Front}, and one of
$F'\subseteq S$ or $F'\cap S=\emptyset$ holds.

Now assume $F\ne\{\emptyset\}$.  By the base disjunction in
\noderef{Front}, together with the definition of the base in
\noderef{FrontBase}, $\mathsf{FrontBase}(F)=X$.  Also $\emptyset\notin F$:
if it were in $F$, then \(\mathsf{InitialSegment}(\emptyset,t)\) would
hold for every $t\in F$ (use the equality case for $t=\emptyset$ and the
least element of $t$ otherwise), so prefix-freeness from
\noderef{Front} would force every member of $F$ to equal $\emptyset$,
contrary to nontriviality.  Thus every member of $F$ is nonempty and has a
least element.

For each $n\in X$, define the ray coloring by
\[
 S_n=\{s\in F_n\mid \{n\}\cup s\in S\}.
\]
By \noderef{FrontRayClosure}, $F_n$ is a front on
$\mathsf{TailSet}(X,n)$.  The preceding tree identification places its ray
rank below that of $F$.  At the diagonal stage below we restrict this ray
before invoking the induction hypothesis.  This is essential: the induction
hypothesis may return a homogeneous sub-front on an arbitrary infinite
subset of the ray's base, so it cannot be required afterwards to lie in a
separately prescribed tail.

We carry out these choices along one increasing sequence, keeping all
infinite-tail facts explicit.  Put $X_{-1}=X$.  Inductively suppose that
$X_{k-1}$ is infinite and $X_{k-1}\subseteq X$.  A nonempty infinite set of
naturals has a least element, so define
\[
 n_k=\min X_{k-1}.
\]
Then $n_k\in X_{k-1}\subseteq X$.  The set
\[
 A_k=X_{k-1}/n_k=\mathsf{TailSet}(X_{k-1},n_k)
\]
where $\mathsf{TailSet}$ is the tail operation of \noderef{TailSet}.  This
set is infinite: its complement in $X_{k-1}$ is contained in the finite set
$\{0,1,\ldots,n_k\}$.  First restrict the ray to this prescribed tail:
\[
 K_k=\mathsf{FrontRestrict}(F_{n_k},A_k).
\]
Since $A_k\subseteq\mathsf{TailSet}(X,n_k)$, the restriction lemma
\noderef{FrontRestriction}, applied to the front $F_{n_k}$ from
\noderef{FrontRayClosure}, shows that $K_k$ is a front on $A_k$.
Every prefix of $K_k$ is a prefix of $F_{n_k}$, so its prefix tree is a
subtree of the ray prefix tree.  Thus
\[
 \mathop{\mathrm{rk}}(K_k)\leq\mathop{\mathrm{rk}}(F_{n_k})
 <\mathop{\mathrm{rk}}(F),
\]
using the ray rank comparison above and the well-founded tree from
\noderef{FrontTreeWellFounded}.  Define the restricted coloring
\[
 T_k=\{s\in K_k\mid\{n_k\}\cup s\in S\}.
\]
The induction hypothesis applies to the front $K_k$ and the subset
$T_k\subseteq K_k$.  It gives a homogeneous front $H_k$ on some infinite
base, with $H_k\subseteq K_k$.  The existential-base direction of
\noderef{SubFrontCharacterization}, now applied to $K_k$ and $H_k$, gives an
infinite set
\[
 X_k\subseteq A_k
\]
such that $H_k=K_k\mathbin{|}X_k$.  Since $X_k\subseteq A_k$ and $K_k$ is
the restriction of $F_{n_k}$ to $A_k$, the definition in
\noderef{FrontRestrict} gives
\[
 K_k\mathbin{|}X_k=F_{n_k}\mathbin{|}X_k.
\]
By the definitions of $T_k$ and $S_{n_k}$, the homogeneous alternative for
$H_k=K_k\mathbin{|}X_k$ therefore becomes either
\[
 F_{n_k}\mathbin{|}X_k\subseteq S_{n_k}
 \quad\text{or}\quad
 (F_{n_k}\mathbin{|}X_k)\cap S_{n_k}=\emptyset.
\]
This completes the induction step: $X_k$ is infinite and
$X_k\subseteq X_{k-1}\subseteq X$.  Moreover every element of $X_k$ is
strictly larger than $n_k$, so $n_{k+1}>n_k$; hence $(n_k)_{k\in\mathbb N}$
is strictly increasing.

Encode the first alternative by $b(k)=\mathsf{true}$ and the second by
$b(k)=\mathsf{false}$.  By
\noderef{BooleanInfinitePigeonhole}, for some $c\in\{\mathsf{true},
\mathsf{false}\}$ the index set
\[
 I=\{k\in\mathbb N\mid b(k)=c\}
\]
is infinite.  Define
\[
 Y=\{n_k\mid k\in I\}.
\]
Every $n_k$ lies in $X$, so $Y\subseteq X$.  The set $Y$ is infinite as
well: strict increase makes $k\mapsto n_k$ injective, and if $Y$ were
finite then its preimage $I$ under this injection would be finite, contrary
to the choice of $I$.

It remains to verify the front and the homogeneous conclusion.  By
\noderef{FrontRestriction}, applied to the front $F$ and the infinite
subset $Y\subseteq X$, the family
\[
 F'=F\mathbin{|}Y
   =\mathsf{FrontRestrict}(F,Y)
\]
is a front on $Y$; the equality is the definition in
\noderef{FrontRestrict}, and $F'\subseteq F$ is immediate.

Take $s\in F'$.  Since $s\in F$ and $\emptyset\notin F$, let
$m=\min s$.  Because $m\in\mathsf{FrontBase}(F)=X$ and $s\subseteq Y$,
$m\in Y$, so $m=n_k$ for some $k\in I$.  Strict increase makes this $k$
unique.  Every other element of $s$ is a member $n_j$ of $Y$ larger than
$n_k$, hence $j>k$ and, by the nested construction of the $X_i$, belongs
to $X_k$.  Therefore
\[
 t=s\setminus\{n_k\}\in F_{n_k}\mathbin{|}X_k:
\]
the ray membership follows from the definition in \noderef{Ray}, because
$s\in F$ and $n_k$ is its least element, and the restriction membership
follows from $t\subseteq X_k$.

If $c=\mathsf{true}$, the first alternative at every $k\in I$ gives
$t\in S_{n_k}$, so the definition of $S_{n_k}$ gives
\[
 s=\{n_k\}\cup t\in S.
\]
Thus $F'\subseteq S$.  If $c=\mathsf{false}$, the second alternative
gives $t\notin S_{n_k}$, and the same definition gives $s\notin S$.
Thus $F'\cap S=\emptyset$.  In either case $F'$ and $Y$ satisfy the
claimed conclusion.
\end{proof}
